\begin{lemma}[Jordan-neighborhood basis for full continua]
\label{lem:jordan-basis}
Let \(K\Subset U\subset\mathbb C\) be a nonempty full continuum and
let \(\mathcal O\) be any neighborhood of \(K\) compactly contained in
\(U\).  Then there is a Jordan domain \(V\) with
\[
 K\subset V\subset\overline V\subset\mathcal O .
\]
Moreover there is a nested sequence of such Jordan domains with
\[
 K\subset V_{r+1},\qquad \overline V_{r+1}\subset V_r,\qquad
 \bigcap_{r\ge1}\overline V_r=K .
\]
\end{lemma}

\begin{proof}
If \(K\) is a singleton, take round disks.  Otherwise, fullness and
connectedness make
\[
 \widehat{\mathbb C}\setminus K
\]
a simply connected domain containing \(\infty\).  Uniformize it by the
exterior of the closed unit disk, normalized at \(\infty\).  The
inverse images of the circles \(|w|=r\), for regular radii
\(r\downarrow1\), are Jordan curves in the complement of \(K\), and
their bounded complementary sides contain \(K\).  Choosing the radii
sufficiently close to \(1\) places the curve together with its closed
bounded side inside \(\mathcal O\), which gives the first assertion;
choosing the radii monotonically gives the nested sequence.  The
intersection of the closures is \(K\), since every point of the
complement of \(K\) has uniformizing modulus strictly greater than
\(1\) and therefore lies outside \(\overline V_r\) for all radii close
enough to \(1\).
\end{proof}